\chapter{Resultados y discusión}
\label{chap:resultados}

Este capítulo presenta primero la caracterización del banco (temporización y repetibilidad) y
después el contraste del modelo con las medidas. El contraste se organiza alrededor de una pregunta
concreta: \emph{¿se comporta el tambor como un cilindro o como una caracola, y en qué medida?} Todas
las figuras y cifras se generan con los programas de \fichero{memoria/figuras/} a partir de los
ficheros de \fichero{cycles/}, sin edición manual.

\section{Campañas de ensayos}
\label{sec:res-campania}

Se han realizado dos campañas, separadas por el rediseño del ciclo y del modelo.

\begin{description}[style=nextline,leftmargin=1.5em]
  \item[Campaña preliminar (14 y 15 de septiembre)] Quince ficheros con el ciclo anterior, que
        arrancaba en el origen y bajaba hasta el final de carrera (A$\rightarrow$B$\rightarrow$A), y
        con el modelo que situaba la salida del cable en la vertical de la articulación. Catorce
        ensayos válidos, 2744 muestras, doce de ellos a \SI{30}{rpm}.
  \item[Campaña del 21 de septiembre] Quince ficheros con el ciclo reorganizado
        (B$\rightarrow$A$\rightarrow$B, sección~\ref{sec:fw-ciclo}) y la calibración provisional ya
        escrita en la cabecera. De ellos, seis son ciclos completos y \textbf{tres} son ensayos
        limpios y mutuamente consistentes (18:06:37, 18:07:09 y 18:08:46), que son los que se usan
        para el análisis cuantitativo. El resto quedaron interrumpidos o con el cable mal asentado,
        y se descartan de forma explícita en la tabla~\ref{tab:ensayos21}.
\end{description}

\begin{table}[H]
\centering
\caption{Ensayos de la campaña del 21 de septiembre con ciclo completo.}
\label{tab:ensayos21}
\small
\begin{tabularx}{\textwidth}{@{}lrrrXl@{}}
\toprule
Ensayo & Muestras & $p_{\max}$ (rev) & $R^2$ fuerza & Observaciones & Uso \\
\midrule
18:06:37 & 180 & 1,789 & 0,932 & ciclo limpio & sí \\
18:07:09 & 180 & 1,806 & 0,931 & ciclo limpio & sí \\
18:08:46 & 180 & 1,810 & 0,930 & ciclo limpio & sí \\
18:16:41 & 172 & 1,847 & 0,187 & la señal de la célula apenas varía & no \\
18:21:20 & 176 & 1,867 & 0,165 & ídem, con par de bajada muy bajo & no \\
18:34:24 & 180 & 1,864 & 0,956 & rango de fuerza distinto: tara cambiada & no \\
\bottomrule
\end{tabularx}
\end{table}

\section{Temporización del lazo de control}
\label{sec:res-temporizacion}

La figura~\ref{fig:temporizacion} muestra la distribución del periodo entre muestras y del tiempo de
trabajo del lazo, medidos por el propio firmware. En la campaña del 21 de septiembre:

\begin{itemize}
  \item el periodo medio es \SI{50.004}{\milli\second} con una desviación típica de
        \SI{0.059}{\milli\second}, y el \SI{99.7}{\percent} de los periodos son exactamente de
        \SI{50}{\milli\second} (el resto, de \SI{51}{\milli\second}, el grano de \codigo{millis()}).
        La frecuencia de muestreo es de \SI{20.0}{\hertz} y cumple RNF1 con holgura;
  \item el tiempo de trabajo está entre \num{10.84} y \SI{11.00}{\milli\second}, con media de
        \SI{10.97}{\milli\second}: el \SI{22}{\percent} del periodo, con \SI{39}{\milli\second} de
        margen. Su dispersión, de apenas \SI{0.2}{\milli\second}, confirma que el lazo está dominado
        por las tres transacciones Modbus.
\end{itemize}

Las cifras coinciden con las de la campaña preliminar, lo que indica que la temporización no depende
del ensayo.

\begin{figure}[tbp]
\centering
\includegraphics[width=\textwidth]{temporizacion.pdf}
\caption{Izquierda: periodo entre muestras (escala logarítmica). Derecha: tiempo de trabajo por
vuelta del lazo.}
\label{fig:temporizacion}
\end{figure}

\section{Ensayo tipo}

La figura~\ref{fig:ensayo-tipo} muestra el ensayo de las 18:08:46. La secuencia es:

\begin{enumerate}
  \item El eje arranca en B y sube hacia A. Durante la primera media revolución (zona sombreada) la
        señal de la célula permanece plana y el par del motor es de unas \SI{3}{\uacc}: el mecanismo
        está recogiendo holgura y todavía no levanta la masa.
  \item A partir de \SI{0.56}{rev} de avance el par crece con rapidez y la señal de la célula
        empieza a moverse: el cable ya trabaja. Desde ahí hasta A el par sigue creciendo.
  \item Al activarse el final A el ciclo invierte la consigna y el eje baja de vuelta a B, con un par
        menor que en la subida por efecto de las pérdidas.
  \item La repetición se cierra al volver a tocar B. La duración total es de \SI{8.9}{\second}.
\end{enumerate}

\begin{figure}[tbp]
\centering
\includegraphics[width=\textwidth]{n_ensayo_tipo.pdf}
\caption{Ensayo tipo de la campaña del 21 de septiembre: avance, velocidad, señal de la célula con
la tensión del modelo superpuesta, y par del motor. En gris, el tramo inicial sin carga.}
\label{fig:ensayo-tipo}
\end{figure}

\section{Repetibilidad}

Los tres ensayos limpios alcanzan el final de carrera A tras \num{1.789}, \num{1.806} y
\SI{1.810}{rev} de avance, con una desviación típica de \SI{11}{mrev}, y vuelven a B con
un avance final de \num{-14}, \num{+5} y \SI{+1}{mrev}. Es decir, el recorrido entre los
dos finales de carrera se repite dentro del \SI{0.6}{\percent}, y el ciclo cierra prácticamente sobre
el punto de partida. El tramo inicial sin carga también es repetible: el par supera las
\SI{10}{\uacc} en \num{0.55}, \num{0.56} y \SI{0.58}{rev}.

La dispersión combina la repetibilidad de los propios finales de carrera con la deriva de la posición
integrada a lo largo de casi dos revoluciones, de modo que es una cota superior de ambas.

\section{Tensión y par a lo largo del recorrido}

La figura~\ref{fig:fuerza-par} superpone las muestras de los tres ensayos limpios en las que el eje
se mueve de forma estable. Se aprecian tres hechos: los ensayos son muy repetibles, la señal de la
célula recorre el mismo camino en subida y bajada (una vez descontado el retardo de su filtro), y el
par es sistemáticamente mayor en la subida que en la bajada, como corresponde a unas pérdidas que se
oponen siempre al movimiento.

\begin{figure}[tbp]
\centering
\includegraphics[width=\textwidth]{n_fuerza_par.pdf}
\caption{Señal de la célula (izquierda) y par del motor (derecha) frente al avance, en los tres
ensayos limpios. En gris, el tramo sin carga.}
\label{fig:fuerza-par}
\end{figure}

Aplicando la separación de la ecuación~\eqref{eq:separacion} sobre intervalos de \SI{0.05}{rev}
(figura~\ref{fig:descomposicion}) se obtiene, en el tramo con carga (de \num{0.68} a
\SI{1.73}{rev}):

\begin{itemize}
  \item \textbf{par gravitatorio} $\tau_g$: crece de \num{20.8} a \SI{32.3}{\uacc}, es decir,
        un \SI{+56}{\percent};
  \item \textbf{pérdidas} $\tau_f$: \SI{5.4}{\uacc} de media, el \SI{19}{\percent} de $\tau_g$, y
        prácticamente constantes a lo largo del recorrido, lo que corresponde a un rozamiento de tipo
        Coulomb;
  \item \textbf{tensión del modelo} $T$: \emph{decrece} de \num{44.2} a \SI{35.5}{\newton}, un
        \SI{-20}{\percent}, porque al subir la barra el brazo del cable mejora más deprisa de lo que
        crece el momento del peso.
\end{itemize}

\begin{figure}[tbp]
\centering
\includegraphics[width=\textwidth]{n_descomposicion.pdf}
\caption{Arriba: señal de la célula en cada sentido, con la tensión del modelo. Centro: par en cada
sentido, componente gravitatoria $\tau_g$ y pérdidas $\tau_f$. Abajo: radio efectivo medido frente a
las dos orientaciones posibles de la caracola.}
\label{fig:descomposicion}
\end{figure}

\section{Contraste con el modelo: cilindro frente a caracola}
\label{sec:res-contraste}

Éste es el núcleo del trabajo experimental. Se comparan tres hipótesis de tambor, todas con la misma
geometría del mecanismo y el mismo recorrido de cable:

\begin{enumerate}
  \item \textbf{cilindro} de radio constante $r = \SI{34.2}{\milli\metre}$, el que paga el mismo
        cable en el mismo giro;
  \item \textbf{caracola configurada}, con el radio decreciendo de \num{45} a \SI{25}{\milli\metre};
  \item \textbf{caracola invertida}, con el radio creciendo de \num{25} a \SI{45}{\milli\metre}.
\end{enumerate}

\subsection{Lo que dice la fuerza: nada concluyente}

La tabla~\ref{tab:contraste} recoge el ajuste de cada hipótesis contra la señal de la célula. Las
tres dan coeficientes de determinación altos, entre \num{0.93} y \num{0.97}. No es que las tres sean
correctas: es que la señal de la célula, sin calibrar, no tiene capacidad para distinguirlas. El
ajuste absorbe en su pendiente tanto la sensibilidad desconocida como el signo
(sección~\ref{sec:mod-ambiguedad}), y lo poco que diferencia a las tres hipótesis en la forma de la
curva queda por debajo del ruido y del retardo del filtro.

\begin{table}[H]
\centering
\caption{Contraste de las tres hipótesis de tambor sobre el tramo con carga. El $R^2$ del par no
distingue el signo; la columna de la pendiente sí: un valor negativo significa que el modelo y la
medida evolucionan en sentidos opuestos.}
\label{tab:contraste}
\small
\begin{tabularx}{\textwidth}{@{}Xccccc@{}}
\toprule
Hipótesis & \multicolumn{2}{c}{Fuerza} & \multicolumn{2}{c}{Par} & Variación de $\tau$ \\
\cmidrule(lr){2-3}\cmidrule(lr){4-5}
 & $R^2$ & pendiente & $R^2$ & pendiente & del modelo \\
\midrule
Cilindro, $r=\SI{34.2}{\milli\metre}$ & 0,947 & $+$ & 0,780 & $-$ & \SI{-22}{\percent} \\
Caracola configurada, 45$\to$25 mm & 0,966 & $+$ & 0,809 & $-$ & \SI{-45}{\percent} \\
Caracola invertida, 25$\to$45 mm & 0,930 & $+$ & 0,673 & $+$ & \SI{+2}{\percent} \\
\midrule
\textbf{Medido} & \multicolumn{4}{c}{} & \textbf{\SI{+56}{\percent}} \\
\bottomrule
\end{tabularx}
\end{table}

\subsection{Lo que dice el par: el tambor no es un cilindro}

El par del accionamiento no necesita calibración, y con él la comparación sí es concluyente. En un
tambor cilíndrico el par es proporcional a la tensión, porque el radio no cambia:
\begin{equation}
  \tau_{\text{cil}}(\varphi) = T(\varphi)\;r_c \quad\Longrightarrow\quad
  \frac{\Delta\tau}{\tau} = \frac{\Delta T}{T} = \SI{-22}{\percent}.
\end{equation}
Es decir, un cilindro obliga al par a \emph{bajar} un \SI{22}{\percent} a lo largo del tramo con
carga. El par medido \emph{sube} un \SI{56}{\percent}. La discrepancia no es de matiz: modelo y
medida evolucionan en sentidos opuestos, y el cociente entre ambos comportamientos es de un factor
\num{2.0} sobre el recorrido.

La conclusión es inmediata: \textbf{el tambor no se comporta como un cilindro}. Y como la única
diferencia entre las dos hipótesis es el radio de enrollamiento, la medida permite estimarlo: el
radio efectivo de la ecuación~\eqref{eq:radio-efectivo} pasa de \num{0.470} a
\SI{0.909}{\uacc\per\newton}, es decir, \textbf{crece un \SI{93}{\percent}} a lo largo del tramo con
carga (figura~\ref{fig:cilindro}).

\begin{figure}[tbp]
\centering
\includegraphics[width=\textwidth]{n_cilindro.pdf}
\caption{Contraste de las tres hipótesis. Izquierda: tensión del modelo. Centro: par normalizado de
cada hipótesis frente al par medido. Derecha: radio de cada hipótesis frente al radio efectivo
deducido de la medida.}
\label{fig:cilindro}
\end{figure}

\subsection{Cuánto compensa, y en qué sentido}

Que el radio crezca es, cualitativamente, lo que pide la compensación: como la tensión baja al subir
la barra, el radio tiene que crecer para que el producto se mantenga. La cuestión es cuánto. Con el
recorrido medido, para que el par fuese exactamente constante bastaría con que el radio creciera un
\SI{24}{\percent}. El radio efectivo crece un \SI{93}{\percent}, casi cuatro veces más.

El resultado, por tanto, tiene dos caras:

\begin{itemize}
  \item \textbf{La caracola actúa, y de forma medible.} Frente al \SI{-22}{\percent} que impondría un
        cilindro, el par real recorre un camino completamente distinto. El mecanismo cambia la
        relación de transmisión a lo largo del recorrido, que es exactamente lo que se quería
        comprobar.
  \item \textbf{Pero no compensa: sobrecompensa.} En lugar de aplanar el par lo empeora en el otro
        sentido, pasando de una caída del \SI{22}{\percent} a una subida del \SI{56}{\percent}. Con
        estos parámetros, el rizado de par sobre el tramo con carga es peor que el de un cilindro
        equivalente (tabla~\ref{tab:perfiles}).
\end{itemize}

Además, el sentido del radio efectivo medido (creciente) es el \emph{contrario} al de la caracola
configurada en el banco (decreciente, de \num{45} a \SI{25}{\milli\metre}), lo que apunta a que la
pieza está descrita al revés en la configuración, montada al revés en la máquina, o ambas cosas
\verificar{medir el perfil real de la caracola y comprobar en qué sentido se enrolla el cable}.

\section{Ajuste de los parámetros del mecanismo}
\label{sec:res-ajuste}

Los parámetros de la tabla~\ref{tab:geometria} no proceden de una medición, sino de un ajuste sobre
los propios ensayos, y conviene ser explícito sobre su alcance. Ajustando sólo contra la señal de la
célula se obtienen los valores que usa la configuración actual, con $R^2 \approx 0{,}93$. Al repetir
el ajuste incluyendo el par, el optimizador se marcha a valores sin sentido físico ---radio final de
\SI{7}{\milli\metre}, desplazamiento de la salida del cable de \SI{200}{\milli\metre}, reducción
\num{3.5}--- aunque mejore formalmente los coeficientes ($R^2$ de \num{0.99} y \num{0.97}).

Esa divergencia es en sí misma un resultado: significa que el modelo, con la parametrización
elegida, no puede reproducir a la vez la fuerza y el par medidos. Faltan datos independientes, y los
que faltan son precisamente los que no se han medido: las cotas del mecanismo y la calibración de la
célula. Mientras no se midan, cualquier juego de parámetros ajustado es una descripción de los datos,
no una verificación del modelo.

\section{Comparación con la campaña preliminar}

La campaña del 14 y 15 de septiembre, analizada con el modelo anterior (salida del cable en la
vertical, tensión proporcional al cable) y con el ciclo A$\rightarrow$B$\rightarrow$A, sugería lo
contrario: en el tramo central del recorrido la tensión variaba un \SI{17}{\percent} y el par
gravitatorio sólo un \SI{6.4}{\percent}, lo que se interpretó como una compensación parcial. Aquel
resultado no se sostiene, por tres motivos que conviene dejar escritos:

\begin{enumerate}
  \item el modelo de tensión era el del caso particular $dx=0$, que no corresponde al montaje;
  \item el sentido del recorrido era el opuesto y el ciclo no empezaba en el extremo del mecanismo,
        de modo que el tramo analizado no es el mismo;
  \item la calibración de la célula que entonces se daba por buena tenía incluso el signo contrario
        al que se deduce ahora, lo que confirma que la fuerza no estaba midiendo lo que se creía.
\end{enumerate}

La lección metodológica es la que recoge la sección~\ref{sec:prob-signo}: un ajuste con buen $R^2$
sobre una única magnitud sin calibrar puede confirmar una hipótesis equivocada.

\section{Limitaciones de las medidas}
\label{sec:res-limitaciones}

\begin{itemize}
  \item \textbf{Célula sin calibrar.} La constante cuentas/newton ha salido distinta en cada intento
        (\num{14.41} en la configuración, \num{20.5} ajustando contra el modelo, \numrange{27}{28} en
        los ensayos limpios) y con signo cambiante entre campañas. Los valores absolutos de fuerza de
        esta memoria son, por tanto, orientativos.
  \item \textbf{Cotas sin medir.} Toda la tensión del modelo depende de $L_2$, $dx$, $L_3$, $L_4$ y
        $\alpha_0$, que son ajustados y no medidos.
  \item \textbf{Escala del par desconocida.} Por eso todos los resultados se expresan en variaciones
        relativas, que no dependen de ella.
  \item \textbf{Tramo sin carga.} La primera media revolución no aporta información y obliga a
        recortar el recorrido analizado.
  \item \textbf{Filtro de la célula.} La media móvil de ocho lecturas introduce un retardo que sólo
        se compensa en primera aproximación al promediar los dos sentidos.
  \item \textbf{Una sola masa y un solo montaje.} No se ha estudiado la dependencia con la carga.
\end{itemize}

\section{Trabajo pendiente para cerrar la validación}

\begin{enumerate}
  \item \textbf{Medir las cotas} del mecanismo ($L_2$, $dx$, $L_3$, $L_4$) y el ángulo de la barra en
        B, e introducirlas en la configuración en lugar de los valores ajustados.
  \item \textbf{Medir el perfil de la caracola}: radios inicial y final, sentido de enrollamiento y
        barrido, que es la magnitud que el análisis deduce de forma indirecta.
  \item \textbf{Calibrar la célula} con masas patrón: escalones crecientes y decrecientes que cubran
        el rango de uso, ajuste por mínimos cuadrados, linealidad, histéresis y repetibilidad, y
        estimación de la incertidumbre con $k=2$.
  \item \textbf{Eliminar la holgura inicial} o, al menos, medirla, para que el recorrido con carga
        sea el recorrido completo.
  \item \textbf{Repetir la campaña} con todo lo anterior y contrastar de nuevo cilindro frente a
        caracola, que es una comparación que el banco ya sabe hacer.
\end{enumerate}
